L'\textit{Sputnik 1} va ser el primer satèl·lit artificial de la història. Consistia en una esfera d'alumini de 58~cm de diàmetre, que allotjava dins seu l'instrumental científic i de transmissions i amb quatre antenes longitudinals adossades a la part exterior. Tenia una massa de 83,6~kg i el seu període orbital era de 96,2 minuts. Actualment, hi ha rèpliques del satèl·lit en diversos museus del món, com la que es mostra en la fotografia.

\figura[0.35]{fig1.png}

\begin{apartat}{1}
Expliqueu raonadament si l'\textit{Sputnik 1} pot ser considerat un satèl·lit geostacionari. Suposant que l'òrbita hagués estat circular, calculeu-ne l'altura sobre la superfície de la Terra.
\end{apartat}

\begin{apartat}{1}
L'\textit{Sputnik 1} va ser llançat a prop de Baikonur, ciutat del Kazakhstan que es troba a uns 45,5° de latitud nord. A aquesta latitud, els objectes en repòs sobre la superfície de la Terra van a una velocitat d'uns 325~m/s a causa de la rotació del planeta. Calculeu l'energia que va caldre subministrar a l'\textit{Sputnik 1} per a situar-lo en la seva òrbita circular.
\end{apartat}

\dades{%
  $G = 6,67\times 10^{-11}\ \mathrm{N\,m^2\,kg^{-2}}$\\
  $M_\mathrm{Terra} = 5,97\times 10^{24}\ \mathrm{kg}$\\
  $R_\mathrm{Terra} = 6\,370\ \mathrm{km}$}
